\section{Proofs and Computable Accounts for the Bellman Bridge}
\label{app:r15bridge}
Throughout this section, $\langle u,v\rangle_d=d^{-1}u^\top v$ and
$\|u\|_d^2=\langle u,u\rangle_d$. Conditioning on fitted parameters
and selection is understood. The initial distribution has finite
support and maximal dispersion $s_0$. The same arguments apply under
any initial distribution with the stated moment and tail bounds.

\subsection{Finite performance difference and the action bridge}
Bounded drift, bounded controls, clipped innovations, and a finite
horizon give finite state moments of every order. For the deterministic
reference $\pi_k$, the finite composition of smooth transition maps has
derivatives bounded by a polynomial in the initial state and innovations.
The running state reward is bounded and smooth, and the terminal reward
is quadratic. Dominated differentiation therefore establishes
\eqref{eq:r15postdecision} and its first two derivatives. This argument
concerns the finite continuation, so no continuous-time PDE regularity
or derivative-discretization estimate is needed.

Let
$Q_k^\pi(z,b)=\ell_k(z,b)+S_k^\pi\{x_b(z)\}$.
The reference recursion is $V_k^\pi(z)=Q_k^\pi(z,\pi_k)$.
For a possibly history-dependent candidate, condition on its history
at $k$, then telescope $V_{k+1}^\pi(Z_{k+1})-V_k^\pi(Z_k)$.
The two terminal functions agree and the initial law is common, giving
\begin{equation}
 \mathcal R_h(a)-\mathcal R_h(\pi)
 =\sum_k\E\{Q_k^\pi(Z_k,a_k)-Q_k^\pi(Z_k,\pi_k)\}.
 \label{eq:r15telescoping}
\end{equation}
Put $e_k=\widehat q_k-q_k^\pi$ and
$v_k=a_k-\pi_k\bm1$. Since $x_{a_k}-x_{\pi_k}=-hv_k$, the
fundamental theorem of calculus implies, pathwise,
\begin{align*}
 &[\widehat Q_k(Z_k,a_k)-\widehat Q_k(Z_k,\pi_k)]
       -[Q_k^\pi(Z_k,a_k)-Q_k^\pi(Z_k,\pi_k)]\\
 &\hspace{25mm}=-h\int_0^1\langle e_k(X_{k,u}),v_k\rangle_d\,du.
\end{align*}
Cauchy--Schwarz with respect to $h$ times counting measure on the grid,
the candidate probability law, and uniform measure on $[0,1]$ proves
Proposition~\ref{prop:r15bellmanbridge}. The sign of the line integral
is immaterial to its absolute bound. Independence of future reference
innovations from the candidate history is used in the recursion.

\subsection{A computable action gap}
The normalized action Hessian of the approximate continuation problem is
\[
 dD_b^2\widehat Q_k(z,b)
 =-A_k\operatorname{diag}(b_i^{-2})
      -\frac{\zeta A_k}{d}\bm1\bm1^\top
      +h^2dD^2\widehat S_k\{x_b(z)\}
 \preceq\Gamma_k I_d.
\]
The second-order integral remainder on the convex action box proves
\eqref{eq:r15actorgap}. Maximizing its scalar quadratic in each
coordinate gives the stated endpoint or projection rule, including
$\Gamma_k=0$. The value is nonnegative because $b=a_k$ is allowed.

For the quadratic conclusion, let $b_k$ maximize $\widehat Q_k$.
When $-\Gamma_k\geq\mu_0h$, strong concavity and the constrained
first-order inequality at $b_k$ give
\[
 \|a_k-b_k\|_d^2\leq\frac{2\delta_k}{\mu_0h},\qquad
 \widehat Q_k(Z_k,b_k)-\widehat Q_k(Z_k,\pi_k)
                  \geq\frac{\mu_0h}{2}\|b_k-\pi_k\bm1\|_d^2.
\]
Use
$\|b_k-\pi_k\bm1\|_d^2\geq
\|a_k-\pi_k\bm1\|_d^2/2-\|a_k-b_k\|_d^2$
and subtract the actor's gap. Summation yields
$\mathcal M_h\geq\mu_0\mathcal A_h/4-2\mathcal D_h$.
Consequently
\[
 \sqrt{\mathcal A_h\mathcal E_h}
 \leq2\sqrt{(\mathcal M_h+2\mathcal D_h)\mathcal E_h/\mu_0}
 \leq\mathcal M_h/2+\mathcal D_h+2\mathcal E_h/\mu_0.
\]
Together with \eqref{eq:r15bridge}, this proves
\eqref{eq:r15quadratic}. If $a_k=b_k$, the stronger relation
$\mathcal M_h\geq\mu_0\mathcal A_h/2$ proves the stated improvement.
No claim of global concavity of the original stochastic control
functional enters this argument.

\subsection{Continuation targets on one global grid}
The following recursion computes an unbiased normalized derivative
target without a matrix-valued state Jacobian. Start at the supplied
postdecision mean $x$, add innovation $k$, and then simulate the
reference states at $j=k+1,\ldots,N$ with the same global grid and
the disclosed $\pi_j$. With
$\gamma_T=2\chi e^{-\rho T}$, set
\begin{equation}
 \begin{aligned}
 p_N&=-\gamma_T\Pi Z_N,\\
 p_j&=B_j Df(Z_j)^\top\bm1
               +\{I+hDf(Z_j)\}^\top p_{j+1},
       \quad j=N-1,\ldots,k+1.
 \end{aligned}
 \label{eq:r15adjoint}
\end{equation}
The output is $p_{k+1}$; for $k=N-1$ it is $p_N$.
There is no derivative of the reference action because the reference
depends only on the date. Every running consumption term is likewise
constant in the starting state. The chain rule proves
\eqref{eq:r15adjoint}, and the dominated-differentiation argument above
shows its expectation is $q_k^\pi(x)$.

The implemented scalar critic has the explicit decomposition, at
$s=t_{k+1}$,
\begin{equation}
 \widehat S_k(x)=-\chi e^{-\rho T}
     \{\|\Pi x\|_d^2+\sigma_I^2(1-1/d)(T-s+h)\}
       +(T-s)\{F_0(s)+F_1(s,x)/d\}.
 \label{eq:r15criticform}
\end{equation}
The two networks have distinct parameters. Thus
\[
\widehat q_k(x)=-\gamma_T\Pi x+(T-s)D_xF_1(s,x):
\]
fitting the time level cannot change the spatial costate. Using one
instead of the exact clipped-normal second moment in the known
state-independent noise term changes neither an action difference nor
the costate. The terminal costate is exact. Fitting on a 64-cell grid
and applying this scalar predictor on a 2,048-cell grid does not make
their continuation problems equal; their difference remains in the
assessed $\mathcal E_h$.

For two hidden tanh layers, write
$F_1(s,x)=w_3^\top\tanh\{W_2\tanh(W_1(s,x)+b_1)+b_2\}+b_3$.
Let $W_{1x}$ contain the state columns of $W_1$. Any verified spectral
upper bounds $l_1\geq\|W_{1x}\|_2$,
$l_2\geq\|W_2\|_2$, $l_3\geq\|w_3\|_2$ give
\begin{align}
 \|\widehat q_k(x)+\gamma_T\Pi x\|_d
       &\leq C_{\theta,k}:=(T-s)l_1l_2l_3/\sqrt d,\notag\\
 dD^2\widehat S_k(x)&\preceq\Lambda_k I_d,
       \qquad \Lambda_k=(T-s)l_3l_1^2(l_2^2+l_2).
 \label{eq:r15networkbounds}
\end{align}
The Hessian bound uses $|\tanh'|\leq1$ and $|\tanh''|\leq1$.
The negative quadratic lift may be omitted in this upper bound; an
absolute normalized Hessian bound is $\gamma_T+\Lambda_k$.
These are computable bounds from the saved parameters, not fitted
derivative maxima.

\subsection{Antithetic cancellation and nonasymptotic ranges}
We give explicit ranges for the statistics used in
\eqref{eq:r15statistics}. All displayed constants are upper endpoints
when implemented. The raw costate comparison uses precisely one bank
with $r$ independent antithetic pairs. The executed choice is $r=2$,
or four paths per bank. Set
\[
 \begin{gathered}
 \beta\geq\kappa\|B\|_2,\quad n_k=N-k-1,\quad
 F_k=(1+h\beta)^{n_k},\quad\tau_k=(N-k)h,\\
 P_k=\beta\sum_{j=k+1}^{N-1}B_j(1+h\beta)^{j-k-1}.
 \end{gathered}
\]
A larger enclosed bound $e^{\beta n_kh}$ may replace $F_k$.
A smaller proved bound on $\|Df(z)^\top\bm1\|_d$ may replace the
leading $\beta$ in $P_k$. Suppose
$\|\Pi(a_k-m_0(t_k)\bm1)\|_d\leq\epsilon$.
Define $s_d=\sqrt{d-1}$ and
\begin{align}
 C_k={}&C_{\theta,k}+P_k+\gamma_TF_k\kappa n_kh\notag\\
 &+\gamma_T(F_k-1)\left\{s_0+(\kappa+\epsilon)T
             +\frac{\sigma_I}{\sqrt d}
                    (\sqrt{t_k}+\sqrt{\tau_k})s_d\right\},\notag\\
 a_k^{\rm tail}={}&\gamma_T(F_k-1)\frac{\sigma_I}{\sqrt d}
                       (\sqrt{t_k}+\sqrt{\tau_k/r}),
 \quad C_E=\max_k C_k,\quad A_E=\max_k a_k^{\rm tail}.
 \label{eq:r15costaterange}
\end{align}
Then the ideal statistic $E$ satisfies
\begin{equation}
 \Pr\{|E|>T(C_E+A_Ev)^2\}\leq3e^{-v^2/2},\qquad v\geq0.
 \label{eq:r15etail}
\end{equation}
For the executed choice $v_*=8$, a clipping threshold and an absolute
expectation allowance are
\begin{equation}
 b_E=T(C_E+A_Ev_*)^2,\qquad
 \tau_E=6TA_E e^{-v_*^2/2}(C_E/v_*+A_E).
 \label{eq:r15eclip}
\end{equation}
The same proof bounds $E^{\rm raw}$ by
$2T(C_Z+A_Zv)^2$, where $C_Z,A_Z$ omit $C_{\theta,k}$ from
\eqref{eq:r15costaterange}. Its expectation allowance is twice the
corresponding value of \eqref{eq:r15eclip}. For $E-E^{\rm raw}$,
one can add these two ranges and allowances; equivalently the direct
bound $T(C_E+A_Ev)^2+2T(C_Z+A_Zv)^2$ uses the same three-event
union. Numerical errors are enclosed separately.

To prove these claims, let $J$ be the derivative of a reference
terminal state with respect to $x$. Then
$\|J\|_2\leq F_k$ and $\|J-I\|_2\leq F_k-1$.
Writing the terminal state as its initial value, the accumulated
production, common deterministic drift, and total innovation gives,
after subtracting $q_0(x)=-\gamma_T\Pi x$, the production-adjoint
term plus
\[
 -\gamma_T\{(J^\top-I)\Pi x
           +J^\top\textstyle\sum_{j>k}h\Pi f(Z_j)
           +(J^\top-I)\Pi\xi+\Pi\xi\}.
\]
The average of the last term is exactly zero within an antithetic bank:
the clipping map is odd. The production norm is at most
$\kappa n_kh$. The first term is bounded by $P_k$.
Common noise disappears from $\Pi\xi$, although it still affects the
states and their Jacobians; the bounds on those terms are global.

For the candidate prefix, let
$H_0=\|\Pi\sum_{j<k}\operatorname{clip}_{z}(G_{j,I})\|_2/\sqrt k$
when $k>0$, with zero contribution at $k=0$. The map defining $H_0$
is 1-Lipschitz on its standard Gaussian input and
$\E H_0\leq s_d$. The average of the analogous future norms from
$r$ independent antithetic base paths has mean at most $s_d$ and
Lipschitz constant $r^{-1/2}$. Gaussian concentration therefore bounds
the failure of
$H_0\leq s_d+v$ and of the two future bounds
$H_l\leq s_d+v/\sqrt r$ by $3e^{-v^2/2}$ in total.
On their intersection,
\[
 \|\Pi X_{k,U}\|_d\leq s_0+(\kappa+\epsilon)T
                  +\sigma_I\sqrt{t_k/d}(s_d+v),
\]
and each bank error has norm at most $C_E+A_Ev$.
Cauchy--Schwarz proves \eqref{eq:r15etail}. Integrating its tail,
using $\int_{v_*}^{\infty}e^{-v^2/2}dv\leq e^{-v_*^2/2}/v_*$,
proves \eqref{eq:r15eclip}. Subtracting two bank residuals proves the
raw and paired-error bounds.

For completeness, the predicted improvement has a substantially
smaller range than the sum of two global value bounds. Let
$\Delta\geq\sup_k\|a_k-\pi_k\bm1\|_d$ and let $l_k,u_k$ enclose
the candidate, reference, and line segments between them. Put
\[
 L_\ell=\max_k\max_{m\in\{l_k,u_k\}}
       |(A_k/h)(m^{-1}-\zeta m)-B_k/h|,
 \quad C_\theta=\max_k C_{\theta,k},
\]
and
\begin{align*}
 C_M&=T\Delta\left[L_\ell+C_\theta+
       \gamma_T\{s_0+(\kappa+\epsilon)T+
                        \sigma_I\sqrt{T/d}\,s_d\}\right],\\
 A_M&=T\Delta\gamma_T\sigma_I\sqrt{T/d}.
\end{align*}
The mean value theorem and the prefix event above imply
$\Pr\{|M|>C_M+A_Mv\}\leq e^{-v^2/2}$.
Thus $b_M=C_M+A_Mv_*$ and
$\tau_M=A_Me^{-v_*^2/2}/v_*$ are valid. Also
$0\leq A\leq T\Delta^2$. If $w_{\mathcal A}$ bounds the normalized
diameter of the action boxes and
$\Gamma_+=\max_k\max(0,\Gamma_k/h)$, then
\[
 D\leq T w_{\mathcal A}
       \{L_\ell+C_\theta+\gamma_T\|\Pi x_{a_K}(Z_K)\|_d\}
                    +T\Gamma_+w_{\mathcal A}^2/2.
\]
This gives $b_D,\tau_D$ from the same affine Gaussian-tail calculation.
All ranges remain valid under mixtures over the disclosed initial
profiles and, using the maximum constants, over fitted streams.

\subsection{Statistics, arithmetic, and payoff transfer}
Conditional independence of the two continuation banks gives
\[
 \E[E\mid K,Z_K,U]
       =T\|\widehat q_K(X_{K,U})-q_K^\pi(X_{K,U})\|_d^2,
\]
and the usual two-replicate identity gives
\[
\E[E^{\rm raw}\mid K,Z_K,U]=
T\E[\|Z_1-q_K^\pi(X_{K,U})\|_d^2\mid K,Z_K,U].
\]
Uniform sampling of $K$ proves all the mean identities in
\eqref{eq:r15statistics}. Equal independent samples from each
declared stream identify the finite uniform mean. Independence across
streams and the nonidentical-variable version of the empirical
Bernstein inequality, \citet[Theorem~11]{maurer2009}, justify pooling.
Dependence among $M,A,E,D$ on one path does not invalidate a union
bound across their confidence statements. Clipping changes each
expectation by at most its proved tail allowance.

The finite-grid identities above concern exact real transition and
reward maps with the disclosed coefficients. A protected numerical
audit must additionally enclose its scalar weights, forward-state
error, neural values and derivatives, adjoint recurrence, and statistic
arithmetic. The controller's rounded actions are held fixed when its
rounded state is compared with the exact internal state, so this state
account does not require a derivative of the actor. For example, if
the action bridge has normalized representation error at most $e_x$
and the normalized critic and true continuation costates have
Lipschitz bounds $L_{\widehat q},L_q$ on the enclosing domain,
Minkowski's inequality gives
\[
 \sqrt{\mathcal E_h^{\rm exact}}
 \leq\sqrt{\mathcal E_h^{\rm represented}}
              +\sqrt T(L_{\widehat q}+L_q)e_x.
\]
A deterministic clipped-noise state enclosure may be used to multiply
these small arithmetic errors. It is not used as the statistical
range in place of \eqref{eq:r15costaterange}. The audit saves the
resulting cushions and evaluates the endpoint with outward rounding.
An ordinary platform derivative without such an account remains a
finite-grid diagnostic, rather than a protected confidence endpoint.

The executed account makes these numerical terms explicit. Let
$C_{\max}$ enclose every coordinate of a candidate prefix followed by
a reference continuation, including its numerical representatives.
A sufficient choice is
\[
 C_{\max}=4\left\{\max_y\|y\|_\infty+M_\epsilon T
          +(\sigma_I+\sigma_C)zT/\sqrt h+1\right\}.
\]
The maximum is over the disclosed initial profiles. This deterministic
bound is used for numerical error only. Put $u=2^{-52}$,
$\gamma_n=nu/(1-nu)$, and $\varepsilon_{\tanh}=2^{-34}$.
The inherited polynomial kernels have uniform tanh and consumption-log
error at most $\varepsilon_{\tanh}$. The row dot-product error is bounded
by $\gamma_{d+2}C_{\max}\max_i\sum_j|B_{ij}|$.
Multiplication by $\kappa$, addition of the drift constants and action,
and the protected state update give the saved one-step bound $\xi_h$.
The source evaluates all its positive terms outward, with
$\gamma_{32}$ for the scalar update and an additional
$16u(\sigma_I+\sigma_C)z\sqrt h$ for represented innovations.
Thus
\[
 e_Z=Ne^{\beta T}\xi_h,\qquad
 e_x=(1+h\beta)e_Z+2\xi_h+
                     \gamma_{16}(C_{\max}+hM_\epsilon)
\]
are sufficient prefix and bridge representation bounds. The last term
also covers the line interpolation and a coupling of the binary uniform
representative to a continuous uniform variable within $2^{-53}$.

Let $b_r=\max_i\|B_{i\cdot}\|_2$ and
$b_1=\max_i\sum_j|B_{ij}|$. A sufficient operator error for the
drift derivative used in the adjoint is
\begin{equation}
 \varepsilon_{Df}=\beta\left\{
 b_r\sqrt d\,e_Z+\gamma_{d+2}C_{\max}b_1
       +2\varepsilon_{\tanh}
       +\gamma_5(1+\varepsilon_{\tanh})^2\right\}.
 \label{eq:r15dfarithmetic}
\end{equation}
The first term uses $|\tanh''|\leq1$ to propagate the state error;
the last two enclose the computed $1-\tanh^2$.
If $\widetilde p_{j+1}$ is a computed adjoint and $e_{j+1}$ encloses
its normalized error, a valid next error is
\begin{align}
 e_j={}&(1+h\beta)e_{j+1}
       +(h\|\widetilde p_{j+1}\|_d+B_j^+)
                                      \varepsilon_{Df}
       +\beta r_{B_j}+\omega_j,\label{eq:r15adjointarithmetic}\\
 \omega_j={}&\gamma_{d+24}\left[
       \|\widetilde p_{j+1}\|_d+
       \kappa\||B|\|_2
       \{h\|\widetilde p_{j+1}\|_d+|\bar B_j|\}
                    (1+4\varepsilon_{\tanh})\right]+10^{-24}.
       \notag
\end{align}
Here $B_j^+$ and $r_{B_j}$ are its scalar upper endpoint and midpoint
radius. The terminal error combines an interval evaluation of
$-\gamma_T\Pi\widetilde Z_N$ with $\gamma_Te_Z$.
The coefficient $\gamma_{d+24}$ covers the products, row dot product,
and addition in
$\widetilde p+\{(h\widetilde p+\bar B_j)
\kappa(1-\widetilde{\tanh}^{\,2})\}B$.
The relative dot-product bound uses $\||B|\|_2$, bounded outward by
its Frobenius norm. The small absolute term dominates gradual-underflow
errors in this fixed $d\leq50$ update: the primitive coefficients and
step are bounded, $|B_{ij}|\leq1$, and even a sum of all its elementary
underflow errors is less than $10^{-24}$. The positive recurrence
\eqref{eq:r15adjointarithmetic} itself is evaluated by interval
operations, so no additional empirical error or terminal correction is
inferred from a comparison of grids. Averaging the four path errors
and enclosing the bank's summation gives its recorded error.

An explicit costate Lipschitz bound, used only to multiply $e_x$, is
obtained by putting $H_f=\kappa\|B\|_2^2$ and
$\mathcal B=\sum_jB_j$. The Euclidean bilinear second derivative of
the reference state flow is at most $H_fTe^{2\beta T}$, while the
Hessian of $\overline{f}$ has norm at most $H_f/d$. The chain rule
therefore gives
\begin{equation}
 L_q=H_fe^{2\beta T}\mathcal B(1+\beta\sqrt d\,T)
       +\gamma_Te^{2\beta T}
                          (1+\sqrt d\,C_{\max}H_fT).
 \label{eq:r15qlipschitz}
\end{equation}
For the neural costate one may use
$L_{\widehat q}=\gamma_T+T l_3l_1^2(l_2^2+l_2)$.
The error in the predicted advantage caused by the prefix replacement
is at most
$T\Delta L_{\widehat q}(1+h\beta)e_Z$.
This uses the difference of the two continuation gradients before
taking an absolute bound. The corresponding actor-gap error is at
most $Tw_{\mathcal A}L_{\widehat q}(1+h\beta)e_Z$.
The squared-error account uses the previously displayed Minkowski
bound with $\sqrt T(L_q+L_{\widehat q})e_x$. The source saves all
these constants and the verified positive-curvature indicator.

Finally, numerical enclosures of sample values can be used more
accurately than a uniform arithmetic allowance. Suppose the exact
clipped sample lies in $[L_i,U_i]$, let $m_i$ be the disclosed binary
midpoint, and let $r_i$ enclose
$\max(m_i-L_i,U_i-m_i)$. Projection onto the subspace orthogonal to
the constant vector is nonexpansive. Consequently the exact sample
mean differs from $\bar m$ by at most $\bar r$, and its sample
standard deviation is at most
\[
 s_m+\sqrt{\frac{n}{n-1}}
                  \left(\frac1n\sum_i r_i^2\right)^{1/2}.
\]
With $\ell=\log(4/\alpha_{\rm event})$, it suffices to add
\[
 \bar r+\sqrt{\frac{2\ell}{n-1}}
                  \left(\frac1n\sum_i r_i^2\right)^{1/2}
\]
to each side of the midpoint empirical Bernstein interval. This
construction encloses the unknown exact sample moments. Its
data-dependent numerical radii do not replace the independently proved
population clipping ranges. For the costate product, for example, bank
and critic norm errors $\eta_1,\eta_2$ give the sample radius
$T\{\eta_1\|\widetilde e_2\|_d+
\eta_2\|\widetilde e_1\|_d+\eta_1\eta_2\}$, in addition to the
interval dot product. All mean and variance corrections are evaluated
outward. The raw costate comparison uses the identical represented
queries; the welfare endpoint additionally includes their
exact-state representation account.

Finally, subtraction of the reference transformed reward gives exactly
\[
 \E X_{a,h}=\mathcal R_h(a)-\mathcal R_h(\pi)-D_{\rm hold}
\]
at the level of expectations, where $X_{a,h}$ is the paired statistic
in Proposition~\ref{prop:r11transfer}. The reference internal drift
uses precisely $h\pi_k=\bar M_k$; the difference from the continuous
schedule is already included in that proposition's scalar quadrature
and reference-state account. Pointwise optimality of $m_0(t)$ in the
consumption-only part implies $D_{\rm hold}\geq0$. Applying
Proposition~\ref{prop:r15bellmanbridge}, the simultaneous endpoints,
and the paired payoff-transfer allowance proves
Theorem~\ref{thm:r15mechanism}. Applying
Proposition~\ref{prop:r15actorgap} on the same event proves its
optional quadratic bound. The nonnegative parts of upper risk and
action endpoints are valid because their true expectations are
nonnegative. The maximum of two valid lower bounds is again valid
on that event.
